\documentclass[conference]{IEEEtran}
\IEEEoverridecommandlockouts
\usepackage[utf8]{inputenc}
\usepackage[T1]{fontenc}
\usepackage[brazil]{babel}
\usepackage{cite}
\usepackage{amsmath,amssymb}
\usepackage{array}
\usepackage{tabularx}
\usepackage{booktabs}
\usepackage{url}
\usepackage{xcolor}
\usepackage{rotating}
\usepackage{float}
\usepackage{tikz}
\usetikzlibrary{positioning,arrows.meta,babel,calc}
\usepackage[edges]{forest}
\usepackage[hidelinks]{hyperref}
\renewcommand{\tablename}{Tabela}
\renewcommand{\figurename}{Figura}
\newcolumntype{Y}{>{\raggedright\arraybackslash}X}
\newcolumntype{P}[1]{>{\raggedright\arraybackslash}p{#1}}
\newcommand{\stEV}{\textsc{ev}}
\newcommand{\stPD}{\textsc{pd}}
\newcommand{\stHO}{\textsc{ho}}
\newcommand{\enab}{\ensuremath{\gg}}
\newcommand{\enabinfo}{\ensuremath{[\,]\!\gg}}
\newcommand{\choice}{\ensuremath{[\,]}}
\newcommand{\interl}{\ensuremath{|||}}
\newcommand{\ordind}{\ensuremath{|\!=\!|}}
\newcommand{\deact}{\ensuremath{[>}}
\newcommand{\susp}{\ensuremath{|>}}
\newcommand{\iter}{\ensuremath{^{*}}}
\newcommand{\ev}[1]{\texttt{#1}}
\newcommand{\cena}[3]{\textbf{#1}\\\textit{#2}\\[-2pt]\rule{\linewidth}{0.3pt}\\#3}
% referências cruzadas para tabelas que vivem nos arquivos de partes/
\makeatletter\newcommand{\fakelabel}[2]{\protected@edef\@currentlabel{#2}\label{#1}}\makeatother

% --------------------------------------------------------------- wireframes
\newcommand{\celular}[3]{% posição, título, conteúdo
  \begin{scope}[shift={#1}]
  \draw[rounded corners=5pt, thick] (0,0) rectangle (3.3,6.2);
  \node[font=\scriptsize\bfseries, anchor=south] at (1.65,6.25) {#2};
  \draw (0.1,5.75) rectangle (3.2,5.95);
  \node[font=\tiny] at (0.85,5.85) {Concluir}; \node[font=\tiny] at (2.45,5.85) {Desistir};
  \draw (0.1,0.1) rectangle (3.2,0.5);
  \foreach \x/\t in {0.48/Início,1.25/Rotinas,2.02/Histórico,2.78/Ajuda}{\node[font=\tiny] at (\x,0.3) {\t};}
  #3
  \end{scope}}
\newcommand{\aparelho}[3]{% y, nome, estado
  \node[font=\tiny, anchor=west] at (0.2,#1+0.1) {#2};
  \node[font=\tiny, anchor=west, gray] at (0.2,#1-0.12) {#3};
  \draw[rounded corners=2.5pt] (2.55,#1-0.12) rectangle (3.05,#1+0.12);
  \fill (2.68,#1) circle (0.08);}
\newcommand{\aparelhoL}[2]{% aparelho ligado: botão à direita, trilho preenchido
  \node[font=\tiny, anchor=west] at (0.2,#1+0.1) {#2};
  \node[font=\tiny, anchor=west, gray] at (0.2,#1-0.12) {Ligada};
  \draw[rounded corners=2.5pt, fill=gray!45] (2.55,#1-0.12) rectangle (3.05,#1+0.12);
  \fill[white, draw=black] (2.92,#1) circle (0.08);}

\begin{document}

\title{Design da Interface Instrumentada, do Sistema de Ajuda e do Acoplamento Interface--Agente do UXAgent no Aplicativo CasaFácil}

\author{\IEEEauthorblockN{Maíra Beatriz de Almeida Lacerda}
\IEEEauthorblockA{\textit{Engenharia de Computação} \\
\textit{CEFET-MG -- Campus V}\\
Divinópolis, Brasil \\
maira@aluno.cefetmg.br}
\and
\IEEEauthorblockN{Maria Eduarda Teixeira Souza}
\IEEEauthorblockA{\textit{Engenharia de Computação} \\
\textit{CEFET-MG -- Campus V}\\
Divinópolis, Brasil \\
maria@aluno.cefetmg.br}
}

\maketitle

\begin{abstract}
Este artigo especifica o design da interface instrumentada, do sistema de ajuda e do acoplamento interface--agente do UXAgent, agente proxy baseado em modelo de linguagem que opera um aplicativo no lugar de uma persona. A persona primária é Lúcia, 68 anos, usuária diária de smartphone e novata em automação residencial; a interface é o CasaFácil, aplicativo móvel fictício que controla quatro aparelhos e as rotinas automáticas criadas por outra pessoa da família. A metamensagem diz, numa frase: você pode ver o que a casa fez sozinha, saber quem fez e desfazer ou impedir a mudança, sem perder o que foi programado. O modelo conceitual é o de um painel de aparelhos com interruptores, um diário de mudanças e regras do tipo se--então, operado por manipulação direta e navegação por abas, sem digitação. O agente observa apenas os signos visíveis da tela atual, em texto literal e com ícones descritos pela forma, e age escolhendo o número de uma das falas do usuário que a MoLIC prevê para aquele estado; a resposta é restrita por saída estruturada. A ajuda tem três camadas sob demanda, a dica local, a Ajuda com dez tópicos organizados por tipo de dúvida e, na versão B, o painel Por quê?, e só chega ao agente depois de uma ação registrada. A principal decisão é oferecer a explicação da ação automática sob demanda, e não sempre exibida: ela custa ao especialista 2,7 segundos e um toque, mas torna a leitura da explicação observável no registro e carrega a hipótese de que explicações sob demanda reduzem os erros de atribuição de causa. Ficam congelados a interface, o texto da ajuda, o tradutor de observação, os prompts, as distribuições dos mecanismos e as tarefas de validação.
\end{abstract}

\begin{IEEEkeywords}
Engenharia semiótica, MoLIC, sistema de ajuda, agentes baseados em LLM, pessoas idosas.
\end{IEEEkeywords}

% =====================================================================
\section{Introdução}

O Alinhamento Rastreável (Fase~II)~\cite{lacerda2026faseII} perguntou se um proxy baseado em LLM, parametrizado, produz registros de tarefas úteis para depurar um protocolo e se suas diferenças de desempenho concordam com observações humanas em tarefas não usadas na calibração. A pergunta desta fase antecede os dados: \textit{que decisões de design da interface instrumentada, do sistema de ajuda e do acoplamento interface--agente permitem gerar esses registros sem dar ao agente informação que a persona não teria?}~\cite{lara2026a3}

Para o UXAgent, a interface é o mundo: tudo o que o agente pode saber chega por signos estáticos, dinâmicos e metalinguísticos, e tudo o que pode fazer é uma fala do usuário prevista no modelo de interação~\cite{barbosa2010}. A fronteira de informação da Fase~II é, portanto, a metamensagem do designer. O mesmo artefato tem três leituras: para a persona, é a conversa com o preposto do designer; para o agente, é o espaço de observação e de ação; para as pesquisadoras, é o esquema de eventos. O design falha quando essas leituras divergem sem declaração. Há uma simetria com o Mágico de Oz~\cite{preece2002}: lá um humano simula o sistema; aqui um modelo simula o usuário, e o protótipo só vale se estiver declarado o que é simulado, por quem e com acesso a quê.

Cada decisão recebe status de evidência (\stEV, com fonte), princípio de design (\stPD) ou hipótese operacional (\stHO). As Seções~\ref{sec:modelo}--\ref{sec:ajuda} seguem os itens (a)--(h) do enunciado; as Seções~\ref{sec:prototipo}--\ref{sec:verif} cobrem interface e protótipo, acoplamento e verificação; os modelos completos estão no Apêndice.

% =====================================================================
\section{Do modelo ao design}\label{sec:modelo}

\subsection{Retomada e reformulação}
Projetar reformula o problema~\cite{lara2026a3}. A Tabela~\ref{tab:tucig} reapresenta a formulação $T \times U \times C \times I \times G$ da Fase~II e registra o que mudou ao projetar. O elenco não mudou: (1)~persona primária Lúcia, representada pelo agente com M1--M6 ligados; (2)~comparador, adultos de 30 a 50 anos sem experiência em configurar automação, representado pela mesma arquitetura com distribuições de usuário ``normal''; (3)~antiusuário, quem configura automações (como a filha), não representado.

\begin{table}[!t]
\caption{$T \times U \times C \times I \times G$: Fase II $\rightarrow$ design}
\label{tab:tucig}
\centering
\scriptsize
\setlength{\tabcolsep}{2.5pt}
\begin{tabularx}{\columnwidth}{@{}P{0.25cm}P{2.6cm}Y@{}}
\toprule
& \textbf{Fase II $\rightarrow$ design} & \textbf{Por quê} \\
\midrule
T & 7 telas $\rightarrow$ 9 telas, 3 diálogos (AL, CE1, CE2), ajuda em 3 camadas & Ação irreversível pedia alerta e desfazer; sair sem salvar pedia captura (Seç.~\ref{sec:molic}) \\
U & Inalterado (faixas e distribuições da Fase II) & O design não justifica mudar a persona \\
C & 4 tarefas; $A_{\min}$ de C1: 6/4 $\rightarrow$ 5/5 & Recontagem pelo CTT e KLM; C1 é de calibração. V1 e V2 intocadas \\
I & 17 $\rightarrow$ 21 tipos de evento, mais 4 derivados & Diálogos, camada de ajuda, método escolhido e etiquetas MAC precisam de evento \\
G & Inalterada (sem dados humanos compatíveis) & Curadoria não mudou \\
\bottomrule
\end{tabularx}
\end{table}

\subsection{Metas de design e especificação de usabilidade}
Das metas de usabilidade de Preece \textit{et al.}~\cite[pp.~14--20]{preece2002}, as prioritárias são \textit{segurança} (não estragar o que a filha configurou), \textit{eficácia} (entender e desfazer), \textit{facilidade de aprender} (novata) e \textit{eficiência}. \textit{Facilidade de lembrar} fica fora: o protocolo tem uma sessão. As metas de experiência da persona (não levar susto, não se sentir incapaz, manter o controle) orientam o design (estado escrito, confirmação e desfazer, ajuda discreta), mas não são medidas pelo agente: autorrelato do agente não é medida humana. A Tabela~\ref{tab:especificacao} converte as metas prioritárias em especificação de usabilidade~\cite[cap.~9]{dix2004} com o dicionário de métricas da Fase~II; os níveis valem para pessoas do perfil na versão~B e são \stHO, pois não há medida atual.

\begin{table}[!t]
\caption{Especificação de usabilidade (perfil, versão B; níveis \stHO)}
\label{tab:especificacao}
\centering
\scriptsize
\setlength{\tabcolsep}{2.5pt}
\begin{tabularx}{\columnwidth}{@{}P{1.35cm}Y P{0.75cm}P{0.75cm}P{0.75cm}@{}}
\toprule
\textbf{Meta} & \textbf{Medida (Fase II)} & \textbf{Pior} & \textbf{Plan.} & \textbf{Melhor} \\
\midrule
Eficácia & $S$ por tarefa (verificador externo) & 0,60 & 0,80 & 1,00 \\
Segurança & $R_{\mathrm{ACT}}$ (exclusão confirmada ou mudança não pedida, por oportunidade) & 0,10 & 0,02 & 0 \\
Utilidade da explicação & $R_{\mathrm{ATR}}(A) - R_{\mathrm{ATR}}(B)$ em C1 & 0 & 0,15 & 0,30 \\
Eficiência & $X = A/A_{\min}$ (KLM: $X=1$) & 2,0 & 1,5 & 1,0 \\
Aprender & $U_{\mathrm{pre}}$: 1.º suporte antes do 1.º erro & 0,3 & 0,6 & 1,0 \\
\bottomrule
\end{tabularx}
\end{table}

\subsection{Restrições}
Design é alcançar metas sob restrições~\cite[cap.~5]{dix2004}. Do enunciado: quatro tarefas, duas de validação congeladas; métricas mínimas $S$, $R_k$, $A$ e $U$; sucesso verificado fora do agente; revisão ética (pessoas idosas, sem CEP até agora). \textit{Restrição do comparador.} Os dados humanos de comparação virão de um piloto próprio no mesmo CasaFácil; por construção, ficam fixados pelo comparador: os objetivos das tarefas; a informação disponível (telas, signos, texto da ajuda); as ações essenciais; os mecanismos de prevenção e recuperação (AL, RA, CE1, CE2 e a PA do horário); e o tamanho e a vizinhança dos alvos, que determinam E-SLP. Ficam livres a aparência visual que não altere alvos nem textos (cores, tipografia, ícones das abas) e a ordem de apresentação dos tópicos da ajuda. A interface do estudo de Tsuchiya \textit{et al.}~\cite{tsuchiya2018}, usada só para conferir o nível de $S$, não é conhecida em detalhe; essa comparação é limitação declarada.

% =====================================================================
\section{Metamensagem, espaço de design e modelo conceitual}\label{sec:meta}

\subsection{Metamensagem}
No modelo de metamensagem de de Souza, conforme Barbosa e Silva~\cite[p.~77]{barbosa2010}:
\begin{quote}\footnotesize
\textit{Este é o meu entendimento, como designer, de quem você é:} uma pessoa de 60 anos ou mais que usa o celular todos os dias para conversar e fazer chamadas de vídeo, mas que não criou as automações da casa, lê com atenção em português, prefere resolver sozinha e tem receio de apagar algo por engano; \textit{do que quer ou precisa fazer:} entender por que um aparelho mudou sem que você tocasse nele e desfazer ou impedir a mudança sem estragar o que foi programado; \textit{de como prefere fazê-lo e por quê:} pelo celular, com poucos toques e sem digitar, conferindo cada resultado e recorrendo à Ajuda quando não reconhece um símbolo, para continuar no controle da própria casa. \textit{Este, portanto, é o sistema que projetei para você, e esta é a forma como pode ou deve utilizá-lo:} no Início você vê seus aparelhos, cada um com o estado escrito, e um aviso quando algo muda sozinho; no Histórico vê quem fez cada mudança e, na versão~B, pode perguntar ``Por quê?'' no próprio aviso; em Rotinas pode pausar, editar ou excluir as regras automáticas, e nenhuma exclusão acontece sem sua confirmação, podendo ainda ser desfeita; a Ajuda explica cada símbolo e cada passo, mas só aparece quando você pede.
\end{quote}
Os ícones sem texto e o rótulo em inglês ``Smart Schedule'' contrariam a metamensagem, mas são mantidos de propósito, por reproduzirem aplicativos comerciais~\cite[pp.~34--35]{tsuchiya2018} e exercitarem L2; a metamensagem compensa com a promessa de ajuda (decisão D12, Apêndice).

\subsection{Espaço de design com dois receptores}
A Tabela~\ref{tab:jak} aplica os elementos de Jakobson adotados pela engenharia semiótica~\cite{barbosa2010}; é nela que se decide a equivalência de canal. A versão integral, com as perguntas do designer, está no caderno de design.

\begin{table*}[!t]
\caption{Espaço de design com dois receptores}
\label{tab:jak}
\centering
\begin{tabularx}{\textwidth}{@{}P{1.2cm}P{2.6cm}YYP{3.6cm}@{}}
\toprule
\textbf{Elemento} & \textbf{Pergunta do designer} & \textbf{Para a persona} & \textbf{Para o agente} & \textbf{Diferença declarada e status} \\
\midrule
Emissor & Que limitações e intenções do designer devem ser comunicadas? & O aplicativo mostra o que a casa faz sozinha e permite desfazer; as rotinas foram criadas por outra pessoa; exclusão exige confirmação & As mesmas intenções, reconstruídas só a partir dos signos de $o_t$; nenhuma explicação adicional do designer no prompt & Nenhuma no conteúdo. O agente recebe também a instrução de formato da resposta (escolher um número da lista), que não comunica intenção do designer da interface. \stPD \\
\addlinespace
Receptor & Que aspectos do usuário, como interpretado pelo designer, devem ser comunicados? & Novata em automação, lê português, receio de errar, usa a Ajuda quando insegura & LLM com a ficha da persona e o estado $x_t$ (crença, notas, rótulos aprendidos); comportamento do perfil vem de M1--M6, não do texto & O LLM conhece convenções de aplicativos pelo pré-treinamento (por exemplo, o ícone de duas barras). M2 compensa em parte, com probabilidade $q_{\mathrm{rot}}$. \stHO \\
\addlinespace
Contexto & Que elementos do contexto o sistema processa, e como? & Casa, à tarde, uma mão, sozinha; o aplicativo processa só os estados dos aparelhos e das rotinas & Sem contexto físico; o contexto chega apenas pela instrução literal da tarefa & Sem tempo, cansaço, pressão nem interrupções; diferença declarada nos limites de uso. \stHO \\
\addlinespace
Código & Qual é a linguagem de interface? & Manipulação direta por toque; texto em português; dois ícones sem texto; um rótulo em inglês & Texto em português que descreve a tela na ordem visual; ícones descritos pela forma; ações numeradas & Tradução visual $\rightarrow$ verbal. Descrever a forma do ícone é tomado como equivalente a vê-lo. \stHO \\
\addlinespace
Canal & Por quais canais a metacomunicação chega? & Tela do smartphone (signos estáticos e dinâmicos); Ajuda, dica local e painel ``Por quê?'' (metalinguísticos), só quando abertos & $o_t$ textual a cada passo; os signos metalinguísticos só entram em $o_t$ depois de uma ação registrada de abrir ajuda & Canal diferente, conteúdo equivalente (regra de equivalência de canal). \stPD \\
\addlinespace
Mensagem & O que se quer contar, e com que efeito? & A metamensagem acima: ``você pode ver e desfazer o que a casa faz sozinha'' & A mesma, inferida pelo agente a partir dos signos & Nenhuma intencional; se o agente a reconstruir melhor que a persona, o efeito vem do LLM, não do design. \stHO \\
\bottomrule
\end{tabularx}

\end{table*}

\subsection{Modelo conceitual}
\textit{O que o sistema faz:} mostra o estado de quatro aparelhos da casa, registra quem mudou cada um e guarda regras automáticas (rotinas) criadas por outra pessoa, que a moradora pode suspender, alterar ou apagar. \textit{Como se comporta:} toda mudança aparece na hora, escrita; nada destrutivo acontece sem confirmação; o que foi feito pode ser desfeito; as rotinas agem sozinhas no horário ou na condição que têm. \textit{Com que se parece:} um painel de interruptores, um diário de mudanças e um fichário de regras ``se\dots então''. O estilo de interação é manipulação direta por toque e navegação por abas, sem digitação nem conversa; a perspectiva comunicativa não obriga a interface conversacional~\cite[p.~230]{barbosa2010}. A Tabela~\ref{tab:tipos} indica o tipo de interação predominante por tarefa~\cite[cap.~2]{preece2002}.

\begin{table}[!t]
\caption{Tipo de interação predominante e metáfora por tarefa}
\label{tab:tipos}
\centering
\scriptsize
\setlength{\tabcolsep}{3pt}
\begin{tabularx}{\columnwidth}{@{}P{0.55cm}P{2.6cm}Y@{}}
\toprule
\textbf{Tarefa} & \textbf{Tipo predominante} & \textbf{Metáfora em jogo} \\
\midrule
C1 & Explorar (achar a causa) e manipular (religar) & Diário de mudanças; interruptor \\
C2 & Navegar e instruir (suspender a regra) & Fichário de regras; pausa \\
V1 & Explorar, navegar e instruir & Diário; fichário; pausa \\
V2 & Explorar, navegar e manipular (ajustar o horário) & Diário; fichário; relógio com $+/-$ \\
\bottomrule
\end{tabularx}
\end{table}

\textit{Metáforas e limites}~\cite[cap.~9]{benyon2010}, \cite{barbosa2010}. O \textit{interruptor} sugere que só a pessoa muda o aparelho; a metáfora quebra exatamente quando uma rotina age, e essa quebra é o problema do estudo, compensado pelo aviso e pelo Histórico. A \textit{rotina} lembra um hábito diário, e não uma regra que age sozinha; por isso a Ajuda~H1 a define como ``regra que o aplicativo segue sozinho''. O \textit{diário} não mostra intenções, só autores; a versão~B acrescenta o motivo. A \textit{pausa} de mídia (duas barras) é convenção que a persona pode não conhecer (L2). \textit{Conceitual e físico}~\cite[p.~199]{benyon2010}: o design conceitual é o que está acima, com os objetos aparelho, rotina e registro de mudança e as operações ver, desfazer, suspender, alterar e apagar; o design físico (Seção~\ref{sec:prototipo}) decide telas, rótulos, ícones e posições.

% =====================================================================
\section{Mapa de objetivos e cenários de interação}\label{sec:cenarios}

\subsection{Mapa de objetivos}
A Tabela~\ref{tab:objetivos} classifica os objetivos~\cite{barbosa2010} da moradora (único papel) sobre as entidades aparelho, rotina e registro. O mapa diz o que a persona quer, nunca como fará.

\begin{table}[!t]
\caption{Mapa de objetivos (papel: moradora)}
\label{tab:objetivos}
\centering
\scriptsize
\setlength{\tabcolsep}{2.5pt}
\begin{tabularx}{\columnwidth}{@{}P{0.45cm}YP{1.05cm}P{0.9cm}@{}}
\toprule
\textbf{Id} & \textbf{Objetivo} & \textbf{Tipo} & \textbf{Entidade} \\
\midrule
F1 & Saber por que um aparelho mudou sem que ela tocasse & Final & registro \\
F2 & Pôr o aparelho no estado que ela quer & Final & aparelho \\
F3 & Impedir que uma regra aja num período sem perdê-la & Final & rotina \\
F4 & Mudar quando uma regra age & Final & rotina \\
I1 & Ver quem fez cada mudança & Instr. direto (F1) & registro \\
I2 & Encontrar a regra que afeta um aparelho & Instr. direto (F3, F4) & rotina \\
I3 & Entender um símbolo ou termo da tela & Instr. indireto & signo \\
I4 & Voltar atrás de um engano & Instr. indireto & todas \\
I5 & Avisar que terminou ou que desistiu & Instr. (protocolo) & tarefa \\
\bottomrule
\end{tabularx}
\end{table}

\subsection{Cenários de interação}
Escritos sem widgets nem rótulos~\cite{barbosa2010}, \cite[cap.~3]{benyon2010}, um por cenário de problema da Fase~II.

\textit{CI1, para CN1 (tarefa crítica), alternativa $\alpha$: explicação sob demanda.} Às 14h02 Lúcia, lendo na sala, nota que a luz apagou. Abre o aplicativo e vê, acima dos aparelhos, o aviso de que a luz da sala foi desligada às 14h02. Pede ao aplicativo o motivo; ele responde que a regra de economia de energia, criada pela filha, apaga a luz da sala às 14h quando há sol. Ela pede para desfazer, vê a luz como ligada, avisa que terminou e diz que foi a regra.

\textit{CI1, alternativa $\beta$: explicação sempre presente.} O próprio aviso já diz que a regra de economia de energia apagou a luz; Lúcia lê, desfaz e conclui, sem precisar pedir o motivo.

\textit{CI1, versão A.} Sem motivo no aviso, Lúcia procura o registro de mudanças; a linha das 14h02 diz que a mudança veio da regra de economia de energia. Ela volta, religa a luz e conclui.

\textit{Perguntas que surgiram e que a análise não levantara.} (1)~O aviso deveria indicar a origem também em A? Não: A é a condição de controle; visibilidade da iniciativa não é explicação (D13). (2)~A explicação deve sumir depois de desfazer? Não: a pessoa confere o resultado. (3)~Desfazer deve depender de abrir a explicação? Não, para que $A_{\min}$ seja igual em A e B. (4)~Em $\beta$, a leitura seria observável? Não, e o uso de suporte ficaria indefinido em B; isso decidiu por $\alpha$ (D4, Seção~\ref{sec:klm}). (5)~A pergunta final sobre a causa deve aceitar ``não sei''? Sim, para não forçar palpites.

\textit{CI2, para CN2.} Na véspera da viagem, Lúcia procura no aplicativo as regras automáticas, reconhece a que liga a cafeteira às 7h, pede que ela fique suspensa e vê que a regra continua guardada, apenas parada. Se pedir para apagá-la, o aplicativo pergunta se tem certeza e, depois, ainda oferece voltar atrás.

\textit{CI3, para CN3 (validação).} A tomada do ventilador desliga de novo; Lúcia pergunta ao aplicativo o motivo, descobre a regra de horário de pico, vai até ela, suspende-a e conclui indicando a regra como causa.

\textit{CI4, para CN4 (validação).} Lúcia descobre que a regra de climatização ligou o ar às 15h, abre a regra, adianta o horário de hora em hora até as 21h, confere o valor, grava e conclui indicando a regra.

CI3 e CI4 só verificam que o design permite realizar V1 e V2; não serviram para escolher entre alternativas. A Tabela~\ref{tab:encadeamento} encadeia cenários, tarefas e objetivos.

\begin{table}[!t]
\caption{Encadeamento: cenário de problema $\rightarrow$ cenário de interação $\rightarrow$ tarefa $\rightarrow$ objetivo}
\label{tab:encadeamento}
\centering
\scriptsize
\begin{tabularx}{\columnwidth}{@{}YYP{1.4cm}P{1.3cm}@{}}
\toprule
\textbf{Problema} & \textbf{Interação} & \textbf{Tarefa} & \textbf{Objetivos} \\
\midrule
CN1 A luz apagou & CI1 ($\alpha$ adotada) & C1 (cal.) & F1, F2 \\
CN2 A viagem & CI2 & C2 (cal.) & F3 \\
CN3 O ventilador parou & CI3 & V1 (val.) & F1, F3 \\
CN4 O ar ligou cedo & CI4 & V2 (val.) & F1, F4 \\
\bottomrule
\end{tabularx}
\end{table}

% =====================================================================
\section{Esquema conceitual de signos}\label{sec:signos}

Os signos são definidos por conteúdo e expressão~\cite{barbosa2010}; o esquema não é o modelo de dados. As Tabelas~\ref{tab:signos1} e~\ref{tab:signos2} preenchem o esquema estendido para os 26 signos usados nas tarefas, com as colunas \textit{Na observação do agente} (regra de equivalência de canal) e \textit{Evento}. A tradução da tela para a observação $o_t$ segue sete regras: (1)~texto literal entre aspas, inclusive em inglês; (2)~ícones descritos pela forma (``ícone de duas barras verticais, sem texto''), nunca pela função; (3)~só a tela atual, sem estado JSON, nomes de variáveis, origem não exibida ou elementos ocultos; (4)~signo dinâmico alterado aparece como ``Mudou na tela: <signo>: <antes> $\rightarrow$ <depois>''; (5)~signos metalinguísticos (ajuda, dica, painel) só após a ação de abri-los; (6)~ações numeradas na ordem visual, com verbo genérico; (7)~nunca HTA, CTT, GOMS, $A_{\min}$, critérios de sucesso ou dados de validação. O CasaFácil não tem botões desabilitados (D2): escrever ``desabilitado'' vazaria estado, e omiti-los diria ao agente que não funcionam. Os signos de aplicação (rotina, situação, Smart Schedule, ícones, ``Por quê?'') têm explicação completa na ajuda~\cite{barbosa2010}.

\begin{table*}[!t]
\caption{Esquema conceitual de signos estendido --- S1, S5, S7, S8 e barras. Origem: dom(ínio), transf(ormado), aplic(ação), conv(encional); tipo: E(stático), D(inâmico), M(etalinguístico)}
\label{tab:signos1}
\centering
\setlength{\tabcolsep}{3pt}
\begin{tabularx}{\textwidth}{@{}P{2.2cm}P{0.9cm}P{2.6cm}P{2.2cm}P{0.5cm}YP{2.3cm}@{}}
\toprule
\textbf{Signo} & \textbf{Origem} & \textbf{Conteúdo (tipo, restrições, default)} & \textbf{Expressão} & \textbf{Tipo} & \textbf{Na observação do agente} & \textbf{Evento} \\
\midrule
aparelho.nome & dom & texto; 4 aparelhos fixos & rótulo à esquerda & E & \texttt{"Luz da sala"} & \ev{screen\_view} \\
aparelho.estado & dom & \{Ligada, Desligada\}; inicial pela tarefa & texto sob o nome + interruptor & D & \texttt{"Desligada"}, interruptor; mudança: ``Mudou na tela: Luz da sala: Desligada $\rightarrow$ Ligada'' & \ev{state\_changed} \\
aviso.mudança & aplic & ``<aparelho>: <estado> às <hh:mm>''; sem origem & faixa no topo de S1 & D & texto literal; nunca a origem & \ev{screen\_view} \\
aviso.porque \textit{[B]} & aplic & botão & texto ``Por quê?'' & E & ``tocar \texttt{"Por quê?"}'' & \ev{support\_open} (por\_que) \\
aviso.desfazer \textit{[B]} & aplic & botão; desfaz a última mudança automática & texto ``Desfazer'' & E & ``tocar \texttt{"Desfazer"}'' & \ev{state\_changed} (origem=pessoa) \\
painel.causa \textit{[B]} & aplic & texto; nome da rotina, ação, hora, condição & painel sob o aviso & M & texto literal, só após tocar ``Por quê?'' & \ev{support\_open}/\ev{close} \\
historico.linha & transf & hora, aparelho, estado, origem \{Você, Rotina <nome>\} & lista em S5 & D & linhas literais, só com S5 aberta & \ev{screen\_view}(S5) \\
barra.navegação & conv & Início, Rotinas, Histórico, Ajuda & abas com ícone e texto & E & ``Barra inferior: \texttt{"Início"}, \texttt{"Rotinas"}, \texttt{"Histórico"}, \texttt{"Ajuda"}'' (o texto basta; ícones das abas omitidos por redundância) & \ev{screen\_view} \\
barra.concluir, barra.desistir & conv (instrumento) & botões; sempre visíveis & texto & E & literal & \ev{screen\_view}(S7); \ev{task\_end} \\
concluir.pergunta & aplic & opções fechadas; nenhuma pré-marcada & botões de opção & E & ``\texttt{"Por que isso aconteceu?"}'' e opções literais; marcada: ``(selecionada)'' porque a tela mostra o círculo preenchido & \ev{cause\_selected} \\
concluir.CE2 & aplic & mensagem de captura & texto vermelho sob a pergunta & M & literal, só depois de Confirmar sem resposta & \ev{dialog\_shown}(CE) \\
\bottomrule
\end{tabularx}

\end{table*}

\begin{table*}[!t]
\caption{Esquema conceitual de signos estendido --- S2, S3, S4, ajuda e diálogos}
\label{tab:signos2}
\centering
\setlength{\tabcolsep}{3pt}
\begin{tabularx}{\textwidth}{@{}P{2.2cm}P{0.9cm}P{2.6cm}P{2.2cm}P{0.5cm}YP{2.3cm}@{}}
\toprule
\textbf{Signo} & \textbf{Origem} & \textbf{Conteúdo (tipo, restrições, default)} & \textbf{Expressão} & \textbf{Tipo} & \textbf{Na observação do agente} & \textbf{Evento} \\
\midrule
rotina.nome & dom (criado pela filha) & texto & item de lista; título em S3 & E & literal & \ev{screen\_view} \\
rotina.situação & aplic & \{Ativa, Pausada\} & texto sob o nome & D & literal; mudança como ``Mudou na tela'' & \ev{state\_changed}(status) \\
rotina.regra & transf & ``Se <condição>, então <ações>'' & texto em S3 & E & literal & \ev{screen\_view}(S3) \\
rotina.pausar/retomar & aplic & alterna a situação & ícone (duas barras $\leftrightarrow$ triângulo), sem texto & E/D & ``botão com ícone de duas barras verticais, sem texto''; com M2: ``[ícone sem nome]'', fora da lista & \ev{action\_executed}; \ev{state\_changed} \\
rotina.editar & aplic & abre S4 & ícone de lápis, sem texto, adjacente ao anterior & E & ``botão com ícone de lápis, sem texto''; M2 como acima & \ev{screen\_view}(S4) \\
rotina.excluir & aplic & ação destrutiva com AL e RA & texto vermelho ``Excluir rotina'' & E & literal (a cor não é descrita: o texto já diz a ação) & \ev{dialog\_shown}(AL) \\
alerta.exclusão (AL) & aplic & confirmação & diálogo modal & M & texto e botões literais & \ev{dialog\_shown}, \ev{dialog\_choice} \\
faixa.desfazer & aplic & RA da exclusão; some ao sair de S2 & faixa em S2 & D & literal enquanto visível & \ev{state\_changed}(restaurada) \\
rotina.horário & aplic & 00h--23h, passo 1~h, sem volta de 23h para 0h (PA); default da rotina & ``Smart Schedule'' + valor + ``$-$''/``$+$'' + ``de 00h a 23h'' & D & literal, em inglês; valor atualizado como ``Mudou na tela'' & \ev{action\_executed} \\
rotina.aparelhos & dom & caixas; marcadas conforme a rotina & caixa + nome & D & ``caixa com marca de seleção, \texttt{"Cafeteira"}'' ou ``caixa vazia, \texttt{"Cafeteira"}'' (a marca é visível) & \ev{action\_executed} \\
editar.salvar/cancelar & conv & botões & texto & E & literal & \ev{state\_changed}; \ev{dialog\_shown}(CE) \\
captura.CE1 & aplic & alterações não salvas & diálogo modal & M & literal & \ev{dialog\_shown}, \ev{dialog\_choice} \\
dica local ``?'' & aplic & abre S9 sobre a tela & botão ``?'' & M & ``tocar \texttt{"?"}''; o balão só após o toque & \ev{support\_open}(dica) \\
ajuda.tópico & aplic & perguntas e respostas (Tabela~\ref{tab:ajuda}) & lista em S6 & M & literal, só após abrir S6 e o tópico & \ev{support\_open}(ajuda, tópico) \\
caderno.nota & --- (artefato da persona) & caminho anotado & papel, fora da tela & --- & ``Sua anotação: ...'' com prob. $r_{\mathrm{nota}}$ (M5) & \ev{note\_shown} (só sintético) \\
\bottomrule
\end{tabularx}

\end{table*}

% =====================================================================
\section{Modelos de tarefa: da HTA ao CTT e ao GOMS}\label{sec:tarefas}

\subsection{CTT}
O CTT~\cite{barbosa2010}, \cite[seç.~18.3]{dix2004} reescreve a HTA das quatro tarefas (Fig.~\ref{fig:ctt-c1} e Tabela~\ref{tab:ctt}, Apêndice) com tarefas abstratas [A], do usuário [U] (as operações não observáveis da Fase~II), interativas [I] (geram eventos de interação) e do sistema [S] (geram eventos do sistema). Em C1, a identificação da causa e o acender da luz são independentes de ordem (\ordind) e habilitam a conclusão com passagem de informação (\enabinfo); as escolhas Histórico/painel e interruptor/Desfazer são \choice. As convenções comuns são:
\begin{equation*}
T = (\mathit{Núcleo}\ \susp\ \mathit{ConsultarAjuda})\ \deact\ \mathit{Encerrar},
\end{equation*}
com a ajuda como suspensão e retomada e a desistência ou o limite de $3A_{\min}$ ações como desativação. Em V1 e V2 o CTT reescreve a HTA congelada sem alterar objetivos, operações ou critérios; três divergências (ordem livre que admite intercalação em C1; \enabinfo{} onde a HTA diz ``sequência''; ajuda global, e não opcional dentro de um subobjetivo) foram registradas, não corrigidas. O CTT não representa erros, tratados na Seção~\ref{sec:molic}.

\subsection{CMN-GOMS de C2}
Na notação de Card, Moran e Newell~\cite{barbosa2010}, \cite[cap.~12]{dix2004}:
\begin{quote}\scriptsize\ttfamily
GOAL: IMPEDIR-QUE-A-CAFETEIRA-LIGUE\\
.\ GOAL: ABRIR-LISTA-DE-ROTINAS: TOCAR Rotinas\\
.\ GOAL: ABRIR-ROTINA: LOCALIZAR "Bom dia"; TOCAR\\
.\ GOAL: SUSPENDER-A-AÇÃO\\
.\ .\ [select METHOD: PAUSAR-ROTINA\\
.\ .\ .\ \ TOCAR ícone de duas barras; VERIFICAR "Pausada"\\
.\ .\ \ \ METHOD: EDITAR-E-DESMARCAR\\
.\ .\ .\ \ TOCAR lápis; TOCAR caixa "Cafeteira";\\
.\ .\ .\ \ TOCAR "Salvar"; VERIFICAR regra]\\
.\ GOAL: ENCERRAR: TOCAR Concluir; TOCAR Confirmar
\end{quote}
Regras de seleção: R1, se outros aparelhos da rotina devem continuar agindo, EDITAR-E-DESMARCAR; R2, senão, PAUSAR-ROTINA. A escolha é observável (\ev{state\_changed} de situação ou de aparelhos) e vira o evento derivado \ev{method\_selected}, indicador e não métrica de sucesso. Regras análogas em C1 geram \ev{method\_selected} \{histórico, painel\} e \{interruptor, desfazer\}.

\subsection{KLM e alternativas da cena principal}\label{sec:klm}
Tabela única de tempos~\cite{card1980}, \cite[cap.~12]{dix2004}: P = 1,10~s, K = 0,28~s, M = 1,35~s; sem H (uma mão) e R = 0 (a latência do LLM é desempenho do sistema). Um toque é P + K = 1,38~s, adaptação ao toque que é \stHO. O operador M segue as mesmas regras em todas as alternativas~\cite[pp.~448--456]{preece2002}: uma M antes do primeiro toque após tela ou painel novo; uma M antes do toque que inicia subobjetivo da HTA; toques repetidos no mesmo botão formam bloco com uma M antes e uma ao final; nenhuma M entre escolher a opção e Confirmar.

A cena principal é o diagnóstico da causa em C1, versão~B. A Tabela~\ref{tab:klm-alt} compara B1 (explicação sob demanda, cenário $\alpha$) e B2 (sempre exibida, $\beta$); a linha de base de todas as tarefas está na Tabela~\ref{tab:klm} (Apêndice).

\begin{table}[!t]
\caption{KLM das alternativas da cena principal (C1, versão B)}
\label{tab:klm-alt}
\centering
\scriptsize
\setlength{\tabcolsep}{2.5pt}
\begin{tabularx}{\columnwidth}{@{}P{1.2cm}YP{0.35cm}P{0.35cm}P{0.65cm}@{}}
\toprule
\textbf{Altern.} & \textbf{Sequência (T = toque)} & \textbf{$n_M$} & \textbf{T} & \textbf{s} \\
\midrule
B1 (adotada) & M T(Por quê?) M T(Desfazer) M T(Concluir) M T(causa) T(Confirmar) & 4 & 5 & 12,30 \\
B2 & M T(Desfazer) M T(Concluir) M T(causa) T(Confirmar) & 3 & 4 & 9,57 \\
\bottomrule
\end{tabularx}
\end{table}

\textit{Decisão D4.} B2 é 2,73~s e um toque mais rápida para o especialista, mas B1 é adotada porque (a)~a pergunta do estudo é sobre explicação sob demanda; (b)~em B1 a leitura é observável (\ev{support\_open(por\_que)}), e em B2 seria operação não observável, deixando $U$ indefinido em B; (c)~B2 adensa uma tela já densa, com mais alvos perto dos interruptores (risco de E-SLP, \stHO). O KLM quantifica o custo aceito; não decide sozinho. B1 carrega a hipótese H-B1: na versão~B, $R_{\mathrm{ATR}}$ é menor que em A, no perfil.

\subsection{Escopo do GOMS e fronteira}
O GOMS descreve especialistas sem erros em tarefas rotineiras~\cite{barbosa2010}, \cite{preece2002}. Aqui serve para comparar alternativas, definir o método de referência e o denominador $A_{\min}$ de $X$, detectar inconsistências (pausar em C2 e em V1 usa o mesmo método) e gerar a ajuda procedimental (H6, H7). Não é modelo da persona: tempos do KLM como tempo de uma novata seriam hipótese sem sustentação, e o UXAgent não simula tempo. HTA, CTT, métodos, regras de seleção e $A_{\min}$ são modelos do analista e nunca chegam ao agente.

% =====================================================================
\section{MoLIC, rupturas e comunicabilidade}\label{sec:molic}

\subsection{Diagrama e leituras}
A MoLIC do único papel de usuário~\cite{barbosa2003molic}, \cite{barbosa2010} está nas Figs.~\ref{fig:molic-a} e~\ref{fig:molic-b} (Apêndice), com cenas, falas de transição, processos do sistema, acesso ubíquo (barras superior e inferior, inclusive Ajuda), falas de recuperação tracejadas e cenas de alerta (AL) e de captura (CE1, CE2) com borda dupla. O vocabulário é controlado (verbo + objeto): ``pausar rotina'' em C2 e em V1 é a mesma fala, o que confirma a consistência entre conversas semelhantes. A Tabela~\ref{tab:leituras} dá as três leituras de cada elemento.

\begin{table*}[!t]
\caption{Leituras de cada elemento da MoLIC ($^*$ só sintético)}
\label{tab:leituras}
\centering
\begin{tabularx}{\textwidth}{@{}P{2.4cm}YYP{3.6cm}@{}}
\toprule
\textbf{Elemento} & \textbf{Leitura da persona} & \textbf{Leitura do agente (observação ou ação)} & \textbf{Evento registrado} \\
\midrule
Cena & Uma tela com um assunto (``ver a casa'', ``ver a rotina'') & Observação $o_t$ com título, signos e lista de ações & \ev{screen\_view} (tela, elementos visíveis) \\
Fala u: & Algo que ela pode tocar ou alternar ali & Item numerado da lista de ações; é a única coisa que o agente pode responder & \ev{action\_proposed}$^*$, \ev{action\_executed} \\
Fala d: & O que a tela mostra depois do toque & Próxima $o'_t$, com a linha ``Mudou na tela'' quando for o caso & \ev{screen\_view} ou \ev{state\_changed} \\
Processo do sistema & O aparelho ou a rotina muda & Invisível como processo; só o resultado aparece na próxima observação & \ev{state\_changed} (variáveis, origem) \\
Fala de recuperação & Desfazer, cancelar, voltar & Ação comum da lista; não é marcada como ``recuperação'' para o agente & \ev{action\_executed}; \ev{dialog\_choice}; etiqueta MAC derivada \\
Alerta e captura & Uma caixa que pergunta ou avisa & Observação da caixa com suas duas opções & \ev{dialog\_shown} (AL ou CE), \ev{dialog\_choice} \\
Acesso ubíquo à ajuda & Aba Ajuda e botão ``?'' em qualquer tela & Itens da lista em toda observação (``tocar Ajuda'', ``tocar ?''); o texto só chega após o toque & \ev{support\_open}, \ev{support\_close} \\
\bottomrule
\end{tabularx}

\end{table*}

\subsection{Máquina de estados}
A MoLIC é ferramenta epistêmica e não modelo formal~\cite{barbosa2010}; a Tabela~\ref{tab:estados} (Apêndice) a traduz numa rede de transição de estados~\cite[cap.~16]{dix2004} implementada no módulo \texttt{casafacil.py}. As falas u: de um estado são exatamente as ações oferecidas ao executor e o conteúdo da observação seguinte; a conclusão de um objetivo é condição do estado avaliada pelo verificador externo; a ajuda é subdiálogo que empilha a cena de origem e devolve a persona a ela.

\subsection{Rupturas}
Para cada erro previsto nas HTAs escolheu-se PP, PA, AL, RA ou CE~\cite{barbosa2010}, verificando as regras ``toda PP tem RA'' e ``toda PA sobre signo de operação vem com PP''. A Tabela~\ref{tab:rup} resume; a matriz completa, com 14 linhas, está no caderno de design. \textit{O mecanismo muda a métrica:} com PA o erro não ocorre (passar de 23h para 0h); com PP e RA ele ocorre e é registrado; com AL a exclusão só conta como E-ACT quando confirmada, e abrir o alerta e cancelar é \ev{dialog\_choice(cancelar)}, etiquetado ``Epa!''. Concluir sem o estado-alvo (E-OMI) não tem PA nem CE: o sistema não conhece o objetivo da tarefa, e avisar entregaria o critério de sucesso. Como pessoas e agente usam o mesmo CasaFácil, as taxas de erro medem a mesma coisa (restrição do comparador).

\begin{table}[!t]
\caption{Rupturas: resumo da matriz (op. da HTA, mecanismo, métrica)}
\label{tab:rup}
\centering
\scriptsize
\setlength{\tabcolsep}{2.5pt}
\begin{tabularx}{\columnwidth}{@{}P{1.6cm}P{1.15cm}YP{0.75cm}@{}}
\toprule
\textbf{Operação} & \textbf{Erro} & \textbf{Mecanismo e evento} & \textbf{Métr.} \\
\midrule
C1.1.4 causa & E-ATR & PP origem em S5 e painel [B]; RA Voltar em S7; \ev{cause\_selected} & $R_{\mathrm{ATR}}$ \\
C1.2.3 interruptor & E-SLP & PP estado escrito; RA tocar de novo; \ev{state\_changed} & $R_{\mathrm{SLP}}$ \\
C2.2.1 excluir & E-ACT & PP texto vermelho afastado; AL; RA faixa Desfazer; \ev{dialog\_shown} & $R_{\mathrm{ACT}}$ \\
C2.2.4 salvar & E-OMI & PP botão Salvar; CE1 ao cancelar; \ev{dialog\_choice} & $R_{\mathrm{OMI}}$ \\
C1.3.3 confirmar & sem resposta & CE2; PP nada pré-marcado & --- \\
V2.3.1 horário & 23h$\rightarrow$0h & PA preso em 00h--23h; PP nota; RA ``$-$'' & --- \\
Concluir & E-OMI & PP instrução; RA Voltar; sem PA/CE & $S$, $R_{\mathrm{OMI}}$ \\
\bottomrule
\end{tabularx}
\end{table}

\subsection{Comunicabilidade operacionalizada}
Oito das treze etiquetas do MAC~\cite{prates2000}, \cite{barbosa2010} são definidas como padrões de eventos (Tabela~\ref{tab:mac}); o verificador grava cada ocorrência como \ev{communicability\_tagged}. Nada é inferido do texto da justificativa do LLM. Cada definição é \stHO: o padrão no registro do agente não equivale à experiência de ruptura de uma pessoa; a mesma definição será aplicada aos registros humanos.

\begin{table*}[!t]
\caption{Etiquetas do MAC como padrões de eventos}
\label{tab:mac}
\centering
\begin{tabularx}{\textwidth}{@{}P{2.2cm}YP{4.4cm}@{}}
\toprule
\textbf{Etiqueta} & \textbf{Assinatura no registro} & \textbf{Onde se espera e o que indica} \\
\midrule
``Cadê?'' & Três ou mais \ev{screen\_view} seguidos por acesso ubíquo ou Voltar, sem \ev{state\_changed} nem \ev{support\_open} entre eles & C1/V1/V2 em A, procurando a origem; C2 procurando as rotinas. Signo de navegação que não comunica onde está a função \\
``O que é isto?'' & \ev{support\_open} na camada dica ou nos tópicos H2, H3 ou H4 & S3 e S4: ícones sem texto e ``Smart Schedule'' (L2) \\
``Epa!'' & \ev{action\_executed} seguido, em até duas ações, da ação inversa sobre o mesmo elemento (alternar de volta, $-$ após $+$, retomar após pausar, Desfazer exclusão) ou \ev{dialog\_choice(cancelar)} no AL & Interruptores adjacentes, ícones adjacentes, $\pm$ (L3, L4) \\
``Onde estou?'' & \ev{screen\_view(X)} seguido de saída de X na ação seguinte, sem nenhuma outra ação em X, duas ou mais vezes na tarefa & Entra e sai de telas sem agir; cena não se identifica \\
``Por que não funciona?'' & A mesma ação executada duas ou mais vezes seguidas sem \ev{state\_changed} (por exemplo, Confirmar sem resposta; $+$ em 23h) & S7 e S4; o sistema não responde como esperado \\
``Socorro!'' & \ev{support\_open} na camada ajuda (S6), qualquer tópico & Qualquer cena; complementa ``O que é isto?'' \\
``Para mim está bom'' & \ev{task\_end(concluida)} e verificador = fracasso & E-OMI e E-ATR: a pessoa acha que terminou \\
``Desisto'' & \ev{task\_end(desistencia)} & Qualquer cena; ruptura total \\
\bottomrule
\end{tabularx}

\end{table*}

% =====================================================================
\section{Sistema de ajuda}\label{sec:ajuda}

A ajuda é a metacomunicação explícita do designer~\cite{barbosa2010}, \cite{silveira2002} e, aqui, variável de protocolo.

\textit{Conteúdo a partir das dúvidas.} Dez tópicos (Tabela~\ref{tab:ajuda}, Apêndice, com o texto exibido) cobrem as dúvidas informativa e motivacional (H1, H10: o que é uma rotina; Concluir e Desistir), descritiva e interpretativa (H2--H4: os dois ícones e ``Smart Schedule''), investigativa e histórica (H5: quem mudou um aparelho), procedimental (H6, H9: suspender; desfazer), de escolha (H7: pausar, editar ou excluir) e navegacional (H8: onde ficam as rotinas). Cada resposta sai de um modelo do analista: esquema de signos (H1--H4), MoLIC (H5, H8), CMN-GOMS (H6, H7) e matriz de rupturas (H7, H9). As respostas são genéricas: nenhuma cita aparelho, rotina ou horário de uma tarefa (D9). A dúvida sugestiva não tem tópico, pois H1, H7 e H8 cobrem as poucas funções; a ausência é decisão.

\textit{Acesso e camadas.} (1)~Dica local ``?'' em S3 e S4, as telas com signos de aplicação sem texto, num balão sobre a própria cena; (2)~Ajuda (S6), aba da barra inferior em todas as cenas, com retorno à cena de origem; (3)~painel ``Por quê?'' (S8, só B), condição experimental registrada como suporte (\texttt{por\_que}) para que $U$ seja comparável entre versões. Os pontos de acesso estão marcados nos esboços da Fig.~\ref{fig:wire}.

\textit{Requisitos de Dix}~\cite[cap.~11]{dix2004}. Disponibilidade: Ajuda em todas as cenas. Precisão e completude: cada signo de aplicação das Tabelas~\ref{tab:signos1}--\ref{tab:signos2} tem tópico, conferido por lista. Consistência: mesmo vocabulário da tela; o ícone é chamado pela forma. Robustez: abrir a ajuda nunca muda aparelhos nem rotinas. Flexibilidade: duas profundidades; não há busca textual porque o protocolo não tem digitação (ausência justificada). Discrição: a ajuda nunca abre sozinha.

\textit{Registro.} Cada acesso gera \ev{support\_open} com camada, tópico, cena de origem e passo, e \ev{support\_close} com a cena de retorno. Uso de suporte por camada: $U_\ell = N(\ev{support\_open}, \ell)/N_{\mathrm{tarefas\ iniciadas}}$; $U = \sum_\ell U_\ell$, com o caderno da persona (só sintético) relatado à parte; abrir tópicos dentro de S6 não conta nova consulta.

\textit{Igualdade e congelamento.} O texto é o mesmo para pessoas e agente, versionado por \textit{hash}, e só chega ao agente quando solicitado (\textit{PedirAjuda})~\cite{lara2026p3}. A ajuda não é adaptativa. A ajuda permitida é a mesma na calibração e na validação; na calibração, nenhum conteúdo foi ajustado pelo desempenho do agente. \textit{Agente de ajuda não é UX-Agent:} o CasaFácil não tem assistente; se tivesse, faria parte do sistema e as lições do Clippy sobre intrusão e controle valeriam para ele~\cite{dix2004}, \cite{preece2002}. O UXAgent está do outro lado e representa a persona.

% =====================================================================
\section{Interface e protótipo}\label{sec:prototipo}

A Fig.~\ref{fig:wire} mostra os esboços das telas centrais: S1 nas versões A, B1 e B2 (as duas alternativas da cena principal) e S3. O inventário completo das nove telas e dos diálogos, na ordem visual que é também a ordem de $o_t$, está na Tabela~\ref{tab:telas} (Apêndice). Decisões de design físico: interruptor acompanhado do estado escrito (D14), para não depender de cor e posição; Concluir e Desistir numa barra fixa (D11), porque o protocolo permite encerrar a qualquer momento e o E-OMI precisa ser possível; ``Excluir rotina'' em texto vermelho, afastado dos ícones adjacentes de pausar e editar (D5).

\begin{figure*}[!t]
\centering
\begin{tikzpicture}[x=1cm,y=1cm]
\celular{(0,0)}{S1 Início -- versão A}{
  \draw[fill=gray!12] (0.15,5.0) rectangle (3.15,5.55);
  \node[font=\tiny, align=left, anchor=west] at (0.2,5.27) {Luz da sala: desligada às 14h02};
  \node[font=\tiny\bfseries, anchor=west] at (0.2,4.7) {Sala e quarto};
  \aparelho{4.25}{Luz da sala}{Desligada}
  \aparelho{3.6}{Ar-condicionado do quarto}{Desligada}
  \aparelho{2.95}{Cafeteira}{Desligada}
  \aparelhoL{2.3}{Tomada do ventilador}
  \node[font=\tiny, gray, align=center] at (1.65,1.3) {origem só no\\Histórico (aba)};}
\celular{(3.8,0)}{S1 -- versão B1 (adotada)}{
  \draw[fill=gray!12] (0.15,4.85) rectangle (3.15,5.55);
  \node[font=\tiny, anchor=west] at (0.2,5.38) {Luz da sala: desligada às 14h02};
  \draw (0.3,4.92) rectangle (1.35,5.15); \node[font=\tiny] at (0.83,5.03) {Por quê?};
  \draw (1.55,4.92) rectangle (2.6,5.15); \node[font=\tiny] at (2.08,5.03) {Desfazer};
  \node[font=\tiny\bfseries, anchor=west] at (0.2,4.55) {Sala e quarto};
  \aparelho{4.1}{Luz da sala}{Desligada}
  \aparelho{3.45}{Ar-condicionado do quarto}{Desligada}
  \aparelho{2.8}{Cafeteira}{Desligada}
  \aparelhoL{2.15}{Tomada do ventilador}
  \node[font=\tiny, gray, align=center] at (1.65,1.2) {causa só no painel S8,\\após ``Por quê?''};}
\celular{(7.6,0)}{S1 -- versão B2 (rejeitada)}{
  \draw[fill=gray!12] (0.15,4.55) rectangle (3.15,5.55);
  \node[font=\tiny, align=left, anchor=west] at (0.2,5.15) {Luz da sala: desligada às 14h02\\pela rotina Economia de energia,\\porque há sol na sala};
  \draw (0.3,4.62) rectangle (1.35,4.85); \node[font=\tiny] at (0.83,4.73) {Desfazer};
  \node[font=\tiny\bfseries, anchor=west] at (0.2,4.25) {Sala e quarto};
  \aparelho{3.8}{Luz da sala}{Desligada}
  \aparelho{3.15}{Ar-condicionado do quarto}{Desligada}
  \aparelho{2.5}{Cafeteira}{Desligada}
  \aparelhoL{1.85}{Tomada do ventilador}
  \node[font=\tiny, gray, align=center] at (1.65,1.05) {leitura não observável;\\tela mais densa};}
\celular{(11.4,0)}{S3 Detalhe da rotina}{
  \node[font=\tiny, anchor=west] at (0.2,5.45) {< Voltar};
  \node[font=\scriptsize\bfseries, anchor=west] at (0.2,5.1) {Bom dia};
  \node[font=\tiny, anchor=west, gray] at (0.2,4.85) {Ativa};
  \node[font=\tiny, align=left, anchor=west] at (0.2,4.4) {Se for 07h00, então ligar\\Cafeteira e Luz da sala};
  \draw (0.3,3.55) rectangle (0.8,3.95); \draw[line width=1.2pt] (0.48,3.63)--(0.48,3.87) (0.62,3.63)--(0.62,3.87);
  \draw (0.9,3.55) rectangle (1.4,3.95); \draw[thick] (1.0,3.62)--(1.3,3.9);
  \draw (2.65,3.55) rectangle (3.05,3.95); \node[font=\tiny] at (2.85,3.75) {?};
  \node[font=\tiny, red!70!black] at (1.65,2.6) {Excluir rotina};
  \node[font=\tiny, gray, align=center] at (1.65,1.4) {ícones sem texto (L2);\\``?'': dica local (S9)};}
\end{tikzpicture}
\caption{Esboços (interface abstrata) de S1 nas versões A, B1 e B2 e de S3. As abas ``Ajuda'' (todas as telas) e ``?'' (S3, S4) são os pontos de acesso à ajuda. B1 e B2 são as alternativas da cena principal (Tabela~\ref{tab:klm-alt}).}
\label{fig:wire}
\end{figure*}

\textit{Protótipo instrumentado.} A interface concreta é a máquina de estados da Tabela~\ref{tab:estados}, implementada em Python (\texttt{casafacil.py}, \texttt{ui\_version = casafacil-1.0}) com estado em JSON, relógio simulado, rotinas reativas e todos os diálogos. A renderização visual para pessoas, sobre a mesma máquina de estados e com os mesmos textos, será construída antes do piloto humano, que depende do CEP; até lá, os esboços acima são a especificação da tela.

% =====================================================================
\section{Acoplamento agente--interface}\label{sec:acoplamento}

O agente segue $a_t \sim \pi_\theta(\cdot \mid o'_t, x_t, \tau)$ e $x_{t+1} = F_\theta(x_t, o_t, a_t)$~\cite{lacerda2026faseII}. A cada passo: (1)~a interface expõe os elementos da tela atual; (2)~o controlador aplica M1 (crença sobre a causa), M2 (ícone não reconhecido vira ``[ícone sem nome]'' e sai da lista) e M5 (notas do caderno) e o tradutor monta $o'_t$; (3)~o LLM recebe só a instrução fixa de formato, a ficha da persona, a instrução literal da tarefa, $o'_t$ e, de $x_t$, a crença, as notas e as últimas oito ações que a persona executou na tarefa, pelo texto visível (memória de trabalho, \stHO), sem histórico de conversa; (4)~a resposta é JSON por saída estruturada, com o campo \texttt{acao} restrito a um \texttt{enum} com os números das falas u: daquele estado; uma resposta inválida gera nova solicitação e depois \texttt{F-FMT}; (5)~o executor aplica M4 (hesitação diante de ação destrutiva) e M3 (deslize para o vizinho de leiaute) e executa; (6)~o registrador grava os eventos. Os parâmetros $\theta$ de M1--M6 são sorteados por sessão a partir de um YAML versionado, com semente. A justificativa do LLM é traço qualitativo, nunca métrica. Exemplo de $o_t$ em S3 (C2, versão~A):

\begin{quote}\scriptsize\ttfamily
Tela: "Bom dia"\quad Situação: "Ativa"\\
Regra: "Se for 07h00, então ligar Cafeteira e Luz da sala"\\
Ações disponíveis: [1] tocar "Concluir" [2] tocar "Desistir" [3] tocar "Voltar" [4] tocar botão com ícone de duas barras verticais, sem texto [5] tocar botão com ícone de lápis, sem texto [6] tocar "?" [7] tocar "Excluir rotina" [8]--[11] abas
\end{quote}

\textit{Diferenças residuais} (\stHO): o agente recebe o texto inteiro da tela a cada passo, e a pessoa pode não ler tudo; a forma descrita de um ícone não é a percepção visual; o LLM pode trazer convenções de aplicativos do pré-treinamento (compensado em parte por M2); o caderno é físico para pessoas e mecanismo registrado no agente.

% =====================================================================
\section{Verificação e congelamento}\label{sec:verif}

\textit{Inspeção de vazamento.} Um teste automático percorre todas as telas de C1 e C2 nas versões A e B e reprova qualquer observação que contenha identificadores internos (nomes de variáveis, de telas ou de rotinas internas), ``HTA'', ``CTT'', ``GOMS'' ou $A_{\min}$; outro confere que ícones aparecem pela forma, e não pela função, e que a origem de uma mudança só aparece no Histórico em A. Testes de importação garantem que o tradutor e a política do LLM não importam os modelos do analista. \textit{Lista de verificação da ajuda:} cada signo de aplicação tem tópico. \textit{Inspeção semiótica:} não realizada (opcional).

\textit{Piloto de instrumentação.} Sessenta sessões sem LLM, só com C1 e C2, versões A e B, perfil e comparador: o executor por regras (só M3 ativo) e uma política aleatória (M1--M6 ativos). Ele verifica o design, não é resultado. Dezoito dos 21 tipos de evento do registro apareceram; os diálogos (AL, CE) e as falhas técnicas não ocorreram e são cobertos por testes. O executor por regras concluiu todas as tarefas ($S = 1$, com $X$ entre 1,00 e 1,06 por causa dos deslizes de M3 e de sua correção), o que confirma que o verificador decide o sucesso só pelo estado final; a política aleatória não concluiu nenhuma ($S = 0$), o que confirma que o verificador não aceita o encerramento declarado. V1 e V2 foram executadas uma única vez pelo executor por regras, sem LLM e sem mecanismos de persona, apenas para conferir que o design permite realizá-las; nada foi ajustado a partir disso, e esses registros não integram nenhum pacote de dados.

\textit{Congelado} antes da calibração com o LLM (Tabela~\ref{tab:congela}). Toda alteração de rótulo, fluxo, ajuda, observação ou prompt depois do início da calibração gera nova versão e recomeça a calibração.

\begin{table}[!t]
\caption{Versões e \textit{hashes} congelados}
\label{tab:congela}
\centering
\scriptsize
\begin{tabularx}{\columnwidth}{@{}P{2.4cm}Y@{}}
\toprule
\textbf{Item} & \textbf{Valor} \\
\midrule
Interface & \texttt{casafacil-1.0} \\
Caderno de design & \texttt{e1-2026-11-22} \\
Texto da ajuda (\textit{hash}) & \texttt{6021a1bf64ae} \\
Mecanismos M1--M6 & \texttt{mech-1.0} \\
Config. perfil / comparador & \texttt{8f32e6a7ef7b} / \texttt{f6833303b58c} \\
Config. genérico / neutro & \texttt{723a1102c3e6} / \texttt{b055335180ff} \\
Modelo & \texttt{claude-sonnet-5-5}, \textit{effort} baixo \\
V1, V2 (instrução, estado, critério) & congeladas em 26/10/2026 \\
\bottomrule
\end{tabularx}
\end{table}

% =====================================================================
\section{Considerações finais}

O design responde à pergunta desta fase com três decisões que se reforçam: a observação do agente contém só os signos que a persona vê (equivalência de canal), suas ações são exatamente as falas u: da MoLIC em cada estado, e todo acesso a explicação ou ajuda é uma ação registrada. A decisão B1 sacrifica eficiência de especialista para tornar a explicação observável, e é a hipótese H-B1 que a calibração e, depois, a validação vão testar. Divergências em relação à Fase~II (\,$A_{\min}$ de C1, proteção da exclusão, conteúdo da observação) estão declaradas, e nenhuma atinge V1 ou V2. A próxima etapa gera os dados sintéticos.

\bibliographystyle{IEEEtran}
\bibliography{referencias}

% =====================================================================
\clearpage
\onecolumn
\renewcommand{\topfraction}{0.98}\renewcommand{\bottomfraction}{0.9}\renewcommand{\textfraction}{0.02}\renewcommand{\floatpagefraction}{0.35}\setcounter{totalnumber}{8}\setcounter{topnumber}{6}\setcounter{bottomnumber}{4}
\appendices
\scriptsize
\section*{Apêndice: modelos completos do design}
CTT e KLM (Fig.~\ref{fig:ctt-c1}, Tabelas~\ref{tab:ctt}--\ref{tab:klm}); MoLIC (Figs.~\ref{fig:molic-a}--\ref{fig:molic-b}); inventário de telas (Tabela~\ref{tab:telas}); conteúdo da ajuda (Tabela~\ref{tab:ajuda}); máquina de estados (Tabela~\ref{tab:estados}); registro de \textit{rationale} (Tabela~\ref{tab:rationale}).


\begin{figure}[!h]
\centering
\resizebox{!}{0.40\textheight}{\resizebox{\textwidth}{!}{%
\begin{forest}
for tree={grow'=east, parent anchor=east, child anchor=west, anchor=west,
  draw, rounded corners=2pt, font=\scriptsize, align=left, inner sep=2pt,
  l sep=4mm, s sep=0.6mm, edge={thin}},
op/.style={draw=none, no edge, font=\footnotesize, inner sep=0pt},
abs/.style={fill=gray!20},
usr/.style={fill=yellow!25},
int/.style={fill=blue!12},
sys/.style={fill=green!15},
[{C1 Restaurar a luz\\e explicar [A]}, abs
  [{Executar C1 [A]}, abs
    [{Núcleo C1 [A]}, abs
      [{Identificar causa [A]}, abs
        [{Pelo Histórico [A]}, abs
          [{Tocar Histórico [I]}, int]
          [\enab, op]
          [{Exibir Histórico [S]}, sys]
          [\enabinfo, op]
          [{Ler origem [U]}, usr]
        ]
        [\choice, op]
        [{Pelo painel \textit{[B]} [A]}, abs
          [{Tocar ``Por quê?'' [I]}, int]
          [\enab, op]
          [{Exibir painel [S]}, sys]
          [\enabinfo, op]
          [{Ler explicação [U]}, usr]
        ]
        [\enabinfo, op]
        [{Decidir causa [U]}, usr]
      ]
      [\ordind, op]
      [{Acender luz [A]}, abs
        [{Pelo interruptor [A]}, abs
          [{[Ir para Início] [I]}, int]
          [\enab, op]
          [{Tocar interruptor [I]}, int]
          [\enab, op]
          [{Atualizar estado [S]}, sys]
        ]
        [\choice, op]
        [{Por Desfazer \textit{[B]} [A]}, abs
          [{Tocar Desfazer [I]}, int]
          [\enab, op]
          [{Atualizar estado [S]}, sys]
        ]
      ]
      [\enabinfo, op]
      [{Concluir [A]}, abs
        [{Tocar Concluir [I]}, int]
        [\enab, op]
        [{Exibir pergunta [S]}, sys]
        [\enab, op]
        [{Escolher causa [I]}, int]
        [\enab, op]
        [{Tocar Confirmar [I]}, int]
        [\enab, op]
        [{Registrar fim [S]}, sys]
      ]
    ]
    [\susp, op]
    [{Consultar ajuda [A]}, abs
      [{Tocar Ajuda ou ``?'' [I]}, int]
      [\enab, op]
      [{Exibir ajuda [S]}, sys]
      [\enabinfo, op]
      [{Ler ajuda [U]}, usr]
      [\enab, op]
      [{Voltar [I]}, int]
    ]
  ]
  [\deact, op]
  [{Encerrar [A]}, abs
    [{Tocar Desistir [I]}, int]
    [\choice, op]
    [{Atingir limite [S]}, sys]
  ]
]
\end{forest}}
}
\caption{Árvore CTT de C1 (calibração), versões A e B. Cinza: abstrata; amarelo: usuário; azul: interativa; verde: sistema.}
\label{fig:ctt-c1}
\end{figure}
\begin{table}[!tp]
\caption{CTT de C2, V1 e V2 em notação textual (V1 e V2 reescrevem a HTA congelada sem alterá-la)}
\label{tab:ctt}
\centering
\begin{tabularx}{\textwidth}{@{}P{1.1cm}Y@{}}
\toprule
\textbf{Tarefa} & \textbf{Expressão CTT} \\
\midrule
C2 & C2 = (Núcleo \susp{} ConsultarAjuda) \deact{} Encerrar \\
& Núcleo = EncontrarRotina \enab{} SuspenderAção \enab{} ConcluirC2 \\
& EncontrarRotina = TocarRotinas[I] \enab{} ExibirLista[S] \enabinfo{} (IdentificarRotina[U] \enabinfo{} TocarRotina[I] \enab{} ExibirDetalhe[S])\iter \quad \textit{(ciclo do plano 1 da HTA)} \\
& SuspenderAção = Pausar \choice{} EditarDesmarcar \\
& Pausar = TocarÍconePausa[I] \enab{} AtualizarSituação[S] \\
& EditarDesmarcar = TocarÍconeEditar[I] \enab{} ExibirEdição[S] \enab{} DesmarcarCafeteira[I] \enab{} TocarSalvar[I] \enab{} SalvarRotina[S] \\
& ConcluirC2 = TocarConcluir[I] \enab{} ExibirConfirmação[S] \enab{} TocarConfirmar[I] \enab{} RegistrarFim[S] \\
\midrule
V1 & V1 = (Núcleo \susp{} ConsultarAjuda) \deact{} Encerrar \\
& Núcleo = IdentificarCausa \enabinfo{} LocalizarRotina \enab{} ImpedirRepetição \enab{} Concluir \quad \textit{(IdentificarCausa e Concluir como em C1)} \\
& LocalizarRotina = TocarIrParaRotina[I]\textit{[B]} \choice{} (TocarRotinas[I] \enab{} ExibirLista[S] \enabinfo{} TocarHorárioDePico[I]) \\
& ImpedirRepetição = Pausar \choice{} EditarDesmarcar(Tomada do ventilador) \\
\midrule
V2 & V2 = (Núcleo \susp{} ConsultarAjuda) \deact{} Encerrar \\
& Núcleo = IdentificarCausa \enabinfo{} AbrirEdição \enab{} AjustarHorário \enab{} TocarSalvar[I] \enab{} SalvarRotina[S] \enab{} Concluir \\
& AbrirEdição = (TocarIrParaRotina[I]\textit{[B]} \choice{} (TocarRotinas[I] \enab{} TocarClimatização[I])) \enab{} TocarÍconeEditar[I] \enab{} ExibirEdição[S] \\
& AjustarHorário = (TocarMais[I] \enab{} ExibirHorário[S] \enabinfo{} ConferirHorário[U])\iter \quad \textit{(até ler 21h00; ciclo do plano 3 da HTA)} \\
\bottomrule
\end{tabularx}

\end{table}
\begin{table}[!tp]
\caption{Linha de base KLM de especialista (caminho mínimo, sem erros). Tempo $= 1{,}35\,n_M + 1{,}38\,A_{\min}$}
\label{tab:klm}
\centering
\begin{tabularx}{\textwidth}{@{}P{1.6cm}YP{0.6cm}P{0.6cm}P{1.0cm}@{}}
\toprule
\textbf{Tarefa / versão} & \textbf{Sequência (M = preparação; T = toque = P+K)} & \textbf{$n_M$} & \textbf{$A_{\min}$} & \textbf{Tempo (s)} \\
\midrule
C1-A & M T(interruptor) M T(Histórico) M T(Concluir) M T(causa) T(Confirmar) & 4 & 5 & 12,30 \\
C1-B & M T(``Por quê?'') M T(Desfazer) M T(Concluir) M T(causa) T(Confirmar) & 4 & 5 & 12,30 \\
C2 (pausar) & M T(Rotinas) M T(``Bom dia'') M T(pausa) M T(Concluir) M T(Confirmar) & 5 & 5 & 13,65 \\
C2 (editar) & M T(Rotinas) M T(``Bom dia'') M T(editar) M T(Cafeteira) T(Salvar) M T(Concluir) M T(Confirmar) & 6 & 7 & 17,76 \\
V1-A & M T(Histórico) M T(Rotinas) M T(``Horário de pico'') M T(pausa) M T(Concluir) M T(causa) T(Confirmar) & 6 & 7 & 17,76 \\
V1-B & M T(``Por quê?'') M T(``Ir para a rotina'') M T(pausa) M T(Concluir) M T(causa) T(Confirmar) & 5 & 6 & 15,03 \\
V2-A & M T(Histórico) M T(Rotinas) M T(``Climatização'') M T(editar) M 6T($+$) M T(Salvar) M T(Concluir) M T(causa) T(Confirmar) & 8 & 14 & 30,12 \\
V2-B & M T(``Por quê?'') M T(``Ir para a rotina'') M T(editar) M 6T($+$) M T(Salvar) M T(Concluir) M T(causa) T(Confirmar) & 7 & 13 & 27,39 \\
\bottomrule
\end{tabularx}

\end{table}


\begin{sidewaysfigure}
\centering
\resizebox{\textheight}{!}{%
\begin{tikzpicture}[
  scene/.style={draw, rounded corners=6pt, text width=3.3cm, align=left, font=\scriptsize, inner sep=3pt, fill=white},
  proc/.style={fill=black, minimum width=6mm, minimum height=3mm, inner sep=0pt},
  ua/.style={draw, fill=gray!25, text width=7.6cm, align=center, font=\scriptsize, inner sep=4pt},
  tr/.style={-{Latex[length=1.8mm]}, thin},
  rec/.style={tr, dashed},
  lab/.style={font=\tiny, align=center, fill=white, inner sep=1pt},
  conn/.style={font=\scriptsize\itshape, align=center}
]
\node[ua] (ua) at (9.2,3.6) {\textbf{Acesso ubíquo}\\barra inferior: Início, Rotinas, Histórico, Ajuda\\barra superior: Concluir, Desistir};
\node[conn] (toS2) at (9.2,5.6) {S2 Rotinas (Fig.~\ref{fig:molic-b})};
\node[circle, fill=black, minimum size=3mm, inner sep=0pt] (open) at (-3.0,0) {};
\node[scene] (s1) at (0,0) {\cena{S1 Início}{Ver a casa}{d: aparelhos e estados\\d: aviso de mudança\\u: alternar aparelho\\u: ver por quê \textit{[B]}\\u: desfazer mudança \textit{[B]}}};
\node[scene] (s8) at (4.8,0) {\cena{S8 Painel ``Por quê?'' [B]}{Entender a ação automática}{d: causa (rotina, hora, condição)\\u: desfazer mudança\\u: ir para a rotina\\u: fechar painel}};
\node[scene] (s5) at (9.6,0) {\cena{S5 Histórico}{Saber quem mudou}{d: hora, aparelho, novo estado, origem}};
\node[scene] (s6) at (14.4,0) {\cena{S6 Ajuda}{Tirar uma dúvida}{d: perguntas (Tab.~\ref{tab:ajuda})\\u: abrir tópico\\u: voltar}};
\node[scene] (s7) at (19.2,0) {\cena{S7 Concluir}{Encerrar a tarefa}{d: Por que isso aconteceu?\\u: escolher resposta\\u: confirmar\\u: voltar}};
\node[proc, label={[font=\tiny]left:alternar aparelho}] (p1) at (-1.0,-3.0) {};
\node[proc, label={[font=\tiny]below:desfazer mudança}] (p2) at (2.4,-3.0) {};
\node[conn] (toS3) at (6.6,-3.0) {S3 Rotina (Fig.~\ref{fig:molic-b})};
\node[conn] (orig6) at (14.4,-3.0) {cena de origem};
\node[proc, label={[font=\tiny]left:registrar fim}] (p8) at (18.4,-3.0) {};
\node[proc, label={[font=\tiny]right:registrar desistência}] (p9) at (27.4,-3.0) {};
\node[conn] (orig7) at (21.6,-3.0) {cena de origem};
\node[scene, double, text width=2.6cm] (ce2) at (24.6,0) {\cena{Captura (CE2)}{Resposta ausente}{d: Escolha uma resposta antes de confirmar.}};
\node[circle, draw, double, double distance=1pt, fill=black, minimum size=3mm, inner sep=0pt] (close) at (22.8,-5.2) {};
\node[font=\tiny, align=left, anchor=west] at (23.1,-5.2) {fechamento:\\verificador externo\\lê o estado};
% abertura
\draw[tr] (open) -- node[lab, above]{d: instrução\\da tarefa} (s1);
% S1
\draw[tr] ([xshift=-6mm]s1.south) -- node[lab, left]{u: alternar\\aparelho} (p1);
\draw[tr] (p1) -- node[lab, right, pos=0.6]{d: novo\\estado} ([xshift=-2mm]s1.south);
\draw[tr] ([xshift=6mm]s1.south) -- node[lab, left]{u: desfazer [B]} (p2);
\draw[tr] ([yshift=2mm]s1.east) -- node[lab, above]{u: ver por quê [B]} ([yshift=2mm]s8.west);
\draw[tr] ([yshift=-4mm]s8.west) -- node[lab, below]{u: fechar} ([yshift=-4mm]s1.east);
\draw[tr] ([xshift=-6mm]s8.south) -- node[lab, right]{u: desfazer} (p2);
\draw[tr] (p2) to[bend left=25] node[lab, below]{d: aparelho restaurado} ([xshift=2mm]s1.south);
\draw[tr] ([xshift=6mm]s8.south) -- node[lab, right]{u: ir para a rotina} (toS3);
% S6 retorno
\draw[tr] (s6) -- node[lab, right]{u: voltar\\d: retorno} (orig6);
% S7
\draw[tr] ([xshift=-4mm]s7.south) -- node[lab, left]{u: confirmar\\(com resposta)} (p8);
\draw[tr] (p8) -- node[lab, left]{d: tarefa encerrada} (close);
\draw[tr] ([yshift=3mm]s7.east) -- node[lab, above]{u: confirmar\\sem resposta} ([yshift=3mm]ce2.west);
\draw[rec] ([yshift=-4mm]ce2.west) -- node[lab, below]{u: escolher\\resposta} ([yshift=-4mm]s7.east);
\draw[tr] ([xshift=12mm]s7.south) -- node[lab, right]{u: voltar} (orig7);
\draw[tr] (p9) -- node[lab, right]{d: tarefa\\interrompida} (close);
% acesso ubíquo
\draw[tr] (ua.south west) -- node[lab, above left]{u: ir para Início} (s1.north);
\draw[tr] ([xshift=4mm]ua.south) -- node[lab, right]{u: ver histórico} (s5.north);
\draw[tr] ([xshift=28mm]ua.south) -- node[lab, right, pos=0.4]{u: pedir ajuda} (s6.north);
\draw[tr] (ua.south east) -- node[lab, above right]{u: concluir} (s7.north);
\draw[tr] (ua.north) -- node[lab, right]{u: ver rotinas} (toS2);
\draw[tr] (ua.east) -| node[lab, above, pos=0.25]{u: desistir} (p9);
\end{tikzpicture}}

\caption{MoLIC, parte (a): navegação, diagnóstico e encerramento. Retângulos arredondados: cenas; pretos: processos do sistema; cinza: acesso ubíquo; tracejado: fala de recuperação de ruptura.}
\label{fig:molic-a}
\end{sidewaysfigure}

\begin{figure}[!tp]
\centering
\resizebox{0.72\textwidth}{!}{\resizebox{\textwidth}{!}{%
\begin{tikzpicture}[
  scene/.style={draw, rounded corners=6pt, text width=3.3cm, align=left, font=\scriptsize, inner sep=3pt, fill=white},
  alert/.style={scene, double},
  proc/.style={fill=black, minimum width=6mm, minimum height=3mm, inner sep=0pt},
  tr/.style={-{Latex[length=1.8mm]}, thin},
  rec/.style={tr, dashed},
  lab/.style={font=\tiny, align=center, fill=white, inner sep=1pt},
  conn/.style={font=\scriptsize\itshape, align=center}
]
\node[conn] (inS2) at (0,3.0) {acesso ubíquo:\\ver rotinas};
\node[conn] (inS3) at (3.4,3.0) {S8: ir para\\a rotina [B]};
\node[scene] (s2) at (0,0) {\cena{S2 Rotinas}{Ver as rotinas}{d: nome e situação\\u: abrir rotina\\u: desfazer exclusão (faixa)}};
\node[scene] (s3) at (5.2,0) {\cena{S3 Rotina}{Ver e controlar a rotina}{d: regra se\dots então, situação\\u: pausar/retomar (ícone)\\u: editar rotina (ícone)\\u: excluir rotina\\u: ver dica (?)\\u: voltar}};
\node[alert] (al) at (10.8,0) {\cena{Alerta (AL)}{Confirmar exclusão}{d: Excluir a rotina X? Ela sai da lista.\\u: excluir\\u: cancelar}};
\node[proc, label={[font=\tiny]left:restaurar rotina}] (p5) at (-1.0,-3.2) {};
\node[proc, label={[font=\tiny]above:pausar/retomar}] (p3) at (7.4,3.0) {};
\node[proc, label={[font=\tiny]right:excluir rotina}] (p4) at (10.8,3.0) {};
\node[scene] (s4) at (5.2,-6.0) {\cena{S4 Editar rotina}{Mudar horário e aparelhos}{d: Smart Schedule hh:00 (00h--23h)\\d: aparelhos com caixas\\u: aumentar/diminuir horário\\u: marcar/desmarcar aparelho\\u: salvar \quad u: cancelar\\u: ver dica (?)}};
\node[proc, label={[font=\tiny]right:salvar rotina}] (p7) at (8.0,-3.0) {};
\node[proc, label={[font=\tiny]right:ajustar valor (PA: 00h--23h)}] (p6) at (5.2,-9.6) {};
\node[alert] (ce) at (11.2,-6.0) {\cena{Captura (CE1)}{Alterações não salvas}{d: Você mudou esta rotina. Salvar as mudanças?\\u: salvar\\u: sair sem salvar}};
\node[scene] (s9) at (0,-6.0) {\cena{S9 Dica local}{Explicar os símbolos da tela}{d: significado dos ícones e de Smart Schedule\\u: fechar}};
\node[conn] (orig9) at (0,-9.6) {cena de origem};
% entradas
\draw[tr] (inS2) -- (s2);
\draw[tr] (inS3) -- ([xshift=-12mm]s3.north);
% S2
\draw[tr] ([yshift=3mm]s2.east) -- node[lab, above]{u: abrir rotina} ([yshift=3mm]s3.west);
\draw[tr] ([yshift=-8mm]s3.west) -- node[lab, below]{u: voltar} ([yshift=-8mm]s2.east);
\draw[rec] ([xshift=-8mm]s2.south) -- node[lab, left]{u: desfazer\\exclusão (RA)} (p5);
\draw[tr] (p5) -- node[lab, right]{d: rotina\\de volta} ([xshift=-1mm]s2.south);
% S3
\draw[tr] ([xshift=10mm]s3.north) -- node[lab, left]{u: pausar|retomar} (p3);
\draw[tr] (p3) -- node[lab, right]{d: situação\\atualizada} ([xshift=17mm]s3.north east);
\draw[tr] ([yshift=4mm]s3.east) -- node[lab, above]{u: excluir rotina} ([yshift=4mm]al.west);
\draw[rec] ([yshift=-6mm]al.west) -- node[lab, below]{u: cancelar} ([yshift=-6mm]s3.east);
\draw[tr] (al) -- node[lab, right]{u: excluir} (p4);
\draw[tr] (p4.north) -- ++(0,0.8) -| node[lab, above, pos=0.25]{d: rotina excluída + faixa ``Desfazer''} ([xshift=-14mm]s2.north);
\draw[tr] ([xshift=-6mm]s3.south) -- node[lab, left]{u: editar rotina} ([xshift=-6mm]s4.north);
\draw[tr] ([xshift=-14mm]s3.south) -- node[lab, above left]{u: ver dica} (s9.north east);
% S4
\draw[tr] ([xshift=10mm]s4.north) -- node[lab, left, pos=0.3]{u: salvar} (p7);
\draw[tr] (p7) -- node[lab, right]{d: rotina salva} ([xshift=12mm]s3.south);
\draw[tr] ([xshift=-4mm]s4.south) -- node[lab, left]{u: aumentar|diminuir|\\marcar|desmarcar} (p6);
\draw[tr] (p6) -- node[lab, right]{d: valor exibido} ([xshift=4mm]s4.south);
\draw[tr] (s4.east) -- node[lab, above]{u: cancelar} node[lab, below]{(com alterações)} (ce.west);
\draw[tr] (ce.north west) -- node[lab, right]{u: salvar} (p7);
\draw[rec] (ce.north) to[out=95, in=-20] node[lab, right, pos=0.5]{u: sair sem salvar} ([xshift=16mm]s3.south east);
\draw[tr] (s4.west) -- node[lab, above]{u: ver dica} (s9.east);
\draw[tr] (s9) -- node[lab, right]{u: fechar\\d: retorno} (orig9);
\end{tikzpicture}}
}
\caption{MoLIC, parte (b): gestão de rotinas. Borda dupla: alerta (AL) e captura de erro (CE). O acesso ubíquo da Fig.~\ref{fig:molic-a} vale em todas as cenas.}
\label{fig:molic-b}
\end{figure}


\begin{table}[!tp]
\caption{Inventário de telas (ordem visual = ordem em $o_t$; [B]: só na versão B)}
\label{tab:telas}
\centering
\begin{tabularx}{\textwidth}{@{}P{2.4cm}Y@{}}
\toprule
\textbf{Tela} & \textbf{Elementos, em ordem} \\
\midrule
Barra superior (todas) & Botão ``Concluir''; botão ``Desistir'' \\
S1 Início & Aviso (se houver mudança automática): ``<Aparelho>: <ligada|desligada> às <hh:mm>'' \textit{[B: + botões ``Por quê?'' e ``Desfazer'']}; título do cômodo; para cada aparelho: nome, texto do estado (``Ligada''/``Desligada''), interruptor. Aparelhos: Luz da sala, Ar-condicionado do quarto, Cafeteira, Tomada do ventilador \\
S2 Rotinas & Título ``Rotinas''; para cada rotina: nome e situação (``Ativa''/``Pausada''); faixa ``Rotina <nome> excluída. Desfazer'' (só logo após uma exclusão, até sair de S2) \\
S3 Detalhe da rotina & Botão ``Voltar''; nome; situação; regra ``Se <condição>, então <ações>''; botão com ícone de duas barras verticais (pausar) ou de triângulo (retomar), sem texto; botão com ícone de lápis (editar), sem texto, adjacente; botão ``?''; botão de texto vermelho ``Excluir rotina'' \\
S4 Editar rotina & Botão ``Cancelar''; ``Smart Schedule'' com o horário (``07h00''), botões ``$-$'' e ``$+$'' adjacentes e a nota ``de 00h a 23h''; ``Aparelhos'': caixa de seleção e nome de cada aparelho; botão ``?''; botão ``Salvar'' \\
S5 Histórico & Título; linhas ``<hh:mm> · <aparelho> · <novo estado> · <origem>'', com origem ``Você'' ou ``Rotina <nome>'', da mais recente à mais antiga \\
S6 Ajuda & Lista de perguntas (Tabela~\ref{tab:ajuda}); ao tocar, o texto da resposta; botão ``Voltar'' \\
S7 Concluir & Pergunta ``Por que isso aconteceu?'' com as opções ``Eu ou outra pessoa'', ``Defeito'', uma por rotina existente, ``Não sei'' (só em C1, V1 e V2); em C2, o texto ``Você terminou?''; botões ``Confirmar'' e ``Voltar'' \\
S8 Painel ``Por quê?'' \textit{[B]} & Texto da causa (``A rotina <nome> <ação> <aparelho> às <hh:mm> porque <condição>''); botões ``Desfazer'', ``Ir para a rotina'' e ``Fechar'' \\
S9 Dica local & Balão sobre S3 ou S4 com o significado dos ícones e campos da tela; botão ``Fechar'' \\
Diálogos & AL: ``Excluir a rotina <nome>? Ela sai da lista de rotinas.'' [Cancelar] [Excluir]. CE1: ``Você mudou esta rotina. Salvar as mudanças?'' [Salvar] [Sair sem salvar]. CE2 (em S7): ``Escolha uma resposta antes de confirmar.'' \\
\bottomrule
\end{tabularx}

\end{table}


\begin{table}[!tp]
\caption{Tópicos da Ajuda (S6) por tipo de dúvida; texto igual ao exibido}
\label{tab:ajuda}
\centering
\begin{tabularx}{\textwidth}{@{}P{0.5cm}P{1.6cm}P{3.0cm}YP{1.8cm}@{}}
\toprule
\textbf{Id} & \textbf{Tipo de dúvida} & \textbf{Pergunta (como aparece)} & \textbf{Resposta} & \textbf{Modelo de origem} \\
\midrule
H1 & Informativa; motivacional & O que é uma rotina? & ``Uma rotina é uma regra que o aplicativo segue sozinho: num horário ou numa situação, ela liga ou desliga aparelhos. Quem cria a rotina escolhe quando ela age. As rotinas ficam na aba Rotinas.'' & Esquema de signos \\
H2 & Descritiva; interpretativa & O que faz o botão de duas barras? & ``É o botão Pausar. Uma rotina pausada fica guardada, mas não age até ser retomada. Quando ela está pausada, o botão mostra um triângulo: é o Retomar.'' & Esquema de signos \\
H3 & Descritiva & O que faz o botão de lápis? & ``É o botão Editar. Ele abre a rotina para mudar o horário e os aparelhos incluídos.'' & Esquema de signos \\
H4 & Interpretativa & O que é ``Smart Schedule''? & ``É o horário em que a rotina age. Toque em $+$ ou $-$ para mudar de hora em hora, de 00h a 23h. Depois, toque em Salvar.'' & Esquema de signos; CTT \\
H5 & Investigativa; histórica & Como saber quem mudou um aparelho? & ``Abra a aba Histórico. Cada linha mostra a hora, o aparelho, o que mudou e quem mudou: `Você' ou o nome da rotina.'' \textit{[B: ``Você também pode tocar em `Por quê?' no aviso da tela Início.'']} & MoLIC (S5, S8) \\
H6 & Procedimental & Como fazer uma rotina parar de agir por um tempo? & ``Abra Rotinas, toque no nome da rotina e toque no botão de duas barras (Pausar). Para ela voltar a agir, toque no triângulo (Retomar).'' & CMN-GOMS (método PAUSAR) \\
H7 & De escolha & Pausar, editar ou excluir: qual usar? & ``Pausar: a rotina para de agir, mas fica guardada. Editar: muda o horário ou os aparelhos. Excluir: apaga a rotina; o aplicativo pede confirmação e, logo depois, ainda mostra `Desfazer' na lista de rotinas.'' & GOMS (regras de seleção); matriz de rupturas \\
H8 & Navegacional & Onde ficam as rotinas? & ``Na aba Rotinas, na barra de baixo, que aparece em todas as telas.'' & MoLIC (acesso ubíquo) \\
H9 & Procedimental (recuperação) & Como desfazer uma mudança? & ``Se você ligou ou desligou um aparelho sem querer, toque no interruptor de novo. Se pausou uma rotina, toque em Retomar. Se excluiu uma rotina, toque em Desfazer na lista de rotinas.'' \textit{[B: ``Se uma rotina mudou um aparelho, toque em Desfazer no aviso.'']} & Matriz de rupturas (RA) \\
H10 & Informativa (protocolo) & O que são Concluir e Desistir? & ``Concluir: avisa que você terminou o que foi pedido. Desistir: encerra sem terminar.'' & Protocolo \\
\bottomrule
\end{tabularx}

\end{table}


\begin{table}[!tp]
\caption{Máquina de estados da interação (UA: acesso ubíquo, disponível em todos os estados exceto diálogos modais)}
\label{tab:estados}
\centering
\begin{tabularx}{\textwidth}{@{}P{1.4cm}P{3.8cm}YP{1.8cm}P{3.2cm}@{}}
\toprule
\textbf{Estado} & \textbf{Falas u: (ações oferecidas)} & \textbf{Processo / fala d:} & \textbf{Próximo} & \textbf{Eventos} \\
\midrule
UA & Início; Rotinas; Histórico; Ajuda; Concluir; Desistir & Ajuda e Concluir empilham a cena de origem; Desistir encerra & S1; S2; S5; S6; S7; FIM & \ev{screen\_view}; \ev{support\_open}; \ev{task\_end(desistencia)} \\
S1 & alternar <aparelho> (4); [B] Por quê?; [B] Desfazer & alternar: novo estado; Desfazer: restaura o aparelho, aviso some & S1; S8; S1 & \ev{action\_executed}, \ev{state\_changed(origem=pessoa)}; \ev{support\_open(por\_que)} \\
S8 [B] & Desfazer; Ir para a rotina; Fechar & Desfazer: restaura; painel mostra ``Feito'' & S8; S3; S1 & \ev{state\_changed}; \ev{screen\_view}; \ev{support\_close} \\
S2 & abrir <rotina>; Desfazer exclusão (se faixa visível) & restaurar rotina excluída & S3; S2 & \ev{screen\_view}; \ev{state\_changed(restaurada)} \\
S3 & Voltar; pausar|retomar (ícone); editar (ícone); ?; Excluir rotina & pausar/retomar: situação alterna & S2; S3; S4; S9; AL & \ev{state\_changed(status)}; \ev{support\_open(dica)}; \ev{dialog\_shown(AL)} \\
AL & Cancelar; Excluir & excluir: rotina sai da lista; faixa ``Desfazer'' em S2 & S3; S2 & \ev{dialog\_choice}; \ev{state\_changed(excluida)} \\
S4 & Cancelar; $-$; $+$; marcar|desmarcar <aparelho>; ?; Salvar & $\pm$: horário $\pm$1~h, preso em 00h--23h (PA); Salvar: grava; Cancelar com alterações: CE1 & S4; S3; CE1; S9 & \ev{action\_executed}; \ev{state\_changed}; \ev{dialog\_shown(CE)} \\
CE1 & Salvar; Sair sem salvar & salvar grava; sair descarta & S3 & \ev{dialog\_choice}; \ev{state\_changed} \\
S5 & (só UA) & --- & --- & \ev{screen\_view(S5)} \\
S6 & abrir <tópico>; Voltar & mostra a resposta; Voltar desempilha & S6; origem & \ev{support\_open(ajuda, tópico)}; \ev{support\_close} \\
S9 & Fechar & desempilha & origem & \ev{support\_close} \\
S7 & escolher <opção>; Confirmar; Voltar & Confirmar com opção: registrar fim; sem opção: CE2 & S7; FIM; S7; origem & \ev{cause\_selected}; \ev{task\_end(concluida)}; \ev{dialog\_shown(CE)} \\
FIM & --- & verificador externo lê o JSON final e decide sucesso & --- & \ev{task\_end}; \ev{error\_tagged}; \ev{method\_selected}; \ev{communicability\_tagged} \\
(controlador) & --- & limite de $3A_{\min}$ ações ou falha técnica & FIM & \ev{task\_end(limite|indeterminado)}; \ev{technical\_failure} \\
\bottomrule
\end{tabularx}

\end{table}


\begin{table}[!tp]
\caption{Decisões de design (Q: questão; O: opções; C: critérios; D: decisão)}
\label{tab:rationale}
\centering
\begin{tabularx}{\textwidth}{@{}P{0.5cm}P{3.0cm}P{3.4cm}YP{0.9cm}@{}}
\toprule
\textbf{Id} & \textbf{Questão} & \textbf{Opções} & \textbf{Critérios e decisão} & \textbf{Status} \\
\midrule
D1 & Como descrever a tela ao agente? & (a) árvore de acessibilidade; (b) lista com identificadores e estados; (c) só signos visíveis em texto literal & Equivalência de canal; sem vazamento. D: (c), regras da Seção~\ref{sec:signos} & \stPD \\
D2 & Usar botões desabilitados? & (a) sim, escrito em $o_t$; (b) sim, omitido da lista; (c) não usar & (a) vaza estado; (b) vaza que não funciona. D: (c), com CE onde seria natural desabilitar & \stPD \\
D3 & Como expressar ícones sem texto ao agente? & (a) função (``pausar''); (b) forma; (c) omitir & (a) vaza significado; (c) esconde o que a persona vê. D: (b), com M2 para a não interpretação & \stHO \\
D4 & Explicação na versão B & B1 sob demanda; B2 sempre exibida & Pergunta do estudo; observabilidade de $U$; densidade da tela; KLM: B2 2,73~s mais rápida. D: B1 & \stPD/\stHO \\
D5 & Proteger ``Excluir rotina'' & (a) sem proteção; (b) AL; (c) AL + RA; (d) PA (remover de S3) & Ação irreversível ao lado de ícones; receio de errar (L4)~\cite[p.~37]{tsuchiya2018}; (d) apagaria a oportunidade de E-ACT. D: (c) & \stEV/\stPD \\
D6 & Sair de S4 sem salvar & (a) descartar calado; (b) salvar automático; (c) CE1 & (b) mudaria a rotina sem intenção explícita; (a) gera E-OMI silencioso. D: (c), justificada por C2, não por V2 & \stPD \\
D7 & Confirmar sem resposta em S7 & (a) desabilitar; (b) pré-marcar ``Não sei''; (c) CE2 & (a) ver D2; (b) enviesa a resposta. D: (c) & \stPD \\
D8 & Camadas da ajuda & (a) só S6; (b) S6 + dica local; (c) ajuda adaptativa & Contextual e em camadas~\cite{barbosa2010}; (c) exigiria regras congeladas e mudaria $U$. D: (b), não adaptativa & \stPD \\
D9 & Conteúdo procedimental da ajuda & (a) genérico, do GOMS; (b) por tarefa & (b) entregaria a solução. D: (a) & \stPD \\
D10 & Limite do horário & (a) circular 23h$\rightarrow$0h; (b) preso em 00h--23h (PA) com nota & (a) gera erro sem sinal. D: (b), PA com PP & \stPD \\
D11 & Onde ficam Concluir e Desistir? & (a) só ao fim; (b) barra fixa & O protocolo permite encerrar a qualquer momento; E-OMI precisa ser possível. D: (b) & \stPD \\
D12 & Manter ícones sem texto e ``Smart Schedule''? & (a) trocar por texto em português; (b) manter & (a) seguiria a metamensagem, mas eliminaria L2, condição do estudo. D: (b), declarada como restrição; compensada pela ajuda & \stEV{} \cite[pp.~34--35]{tsuchiya2018} \\
D13 & Aviso de mudança em A e B & (a) só em B; (b) igual nas duas & Não confundir visibilidade da iniciativa com explicação~\cite{li2026}. D: (b) & \stPD \\
D14 & Estado do aparelho & (a) só interruptor; (b) interruptor + texto & Não depender de cor ou posição; texto literal em $o_t$ sem vazamento. D: (b) & \stPD \\
\bottomrule
\end{tabularx}

\end{table}

\end{document}
